\section{Accessible categories and the relative classifying topos (Part I)}\label{sec:categoryI}

\subsection{The indexing convention and the actual equivalence}
We fix a universe and use ``small'' relative to it.  A category is
$\kappa$-filtered, for an infinite regular cardinal $\kappa$, when every
diagram having fewer than $\kappa$ arrows has a cocone.  An object $c$ is
$\kappa$-presentable when $\operatorname{Hom}(c,-)$ preserves small
$\kappa$-filtered colimits.  The manuscript's ordered sets ``grands devant
$\pi$'' have upper bounds for subsets of cardinal at most $\pi$; their
modern index is therefore $\kappa=\pi^+$, not $\pi$.
For a small category $C$, $\operatorname{Ind}_\kappa(C)$ is the full
subcategory of $[C^{\mathrm{op}},\mathbf{Set}]$ consisting of small
$\kappa$-filtered colimits of representables.
A full subcategory $C\subset E$ is dense when the restricted Yoneda
functor $N_C:E\to[C^{\mathrm{op}},\mathbf{Set}]$ is fully faithful.

\begin{theorem}[Proved: the corrected equivalence]\label{cat:I1}
For a locally small category $E$ with small $\kappa$-filtered colimits,
the following conditions are equivalent.
\begin{enumerate}
\item There is a small full subcategory of $\kappa$-presentable objects
such that every object of $E$ is a $\kappa$-filtered colimit of its objects.
\item $E\simeq\operatorname{Ind}_\kappa(C)$ for a small category $C$.
\item The $\kappa$-presentable objects form an essentially small dense
subcategory $E_\kappa$, and $E_\kappa/X$ is $\kappa$-filtered for every
$X\in E$.
\end{enumerate}
In $\operatorname{Ind}_\kappa(C)$ the $\kappa$-presentable objects are
exactly the retracts of representables.  If $E$ is cocomplete, condition
(3) may omit the last, slice-filteredness condition.  In that case it
also suffices that $E$ have a small dense subcategory of
$\kappa$-presentable objects.
\end{theorem}
\begin{proof}
Representables are $\kappa$-presentable because evaluation commutes with
colimits of presheaves.  If $X=\operatorname{colim}_I yc_i$ is itself
$\kappa$-presentable, its identity factors through some $yc_i$, so it is
a retract.  Conversely a retract of a presentable object is presentable:
the associated Hom functor is a retract of a functor preserving the
colimits in question.  Such retracts do belong to
$\operatorname{Ind}_\kappa(C)$.  For an idempotent $e$ on $yc$, take the
constant-object diagram on the ordinal $\kappa$ whose maps between
strictly increasing indices are $e$; its colimit is the image of $e$.
The ordinal is $\kappa$-filtered by regularity.  There is only a set of
idempotents on representables, which proves essential smallness.

For any presheaf $X$ which is a $\kappa$-filtered colimit of
representables, $C/X$ is $\kappa$-filtered.  Indeed, every family of
fewer than $\kappa$ elements, and all relations in a diagram between
them, is represented at a common stage, and the fewer than $\kappa$
required equalities hold at a further common stage.  That stage defines
a cocone in $C/X$.  The usual presheaf presentation
$X=\operatorname{colim}_{C/X}yc$ therefore gives the desired canonical
filtered presentation.  The same element argument shows that a
$\kappa$-filtered colimit of such presheaves again has a
$\kappa$-filtered category of elements.  Thus these colimits remain in
$\operatorname{Ind}_\kappa(C)$ and are computed pointwise.

Suppose (1) holds, with the chosen small category $C$.  For a
presentation $X=\operatorname{colim}_I c_i$, presentability gives
$N_CX=\operatorname{colim}_I yc_i$.  Hence, for $Y\in E$,
\[
 \operatorname{Hom}(N_CX,N_CY)
 =\lim_i\operatorname{Hom}(c_i,Y)
 =\operatorname{Hom}_E(X,Y).
\]
The comparison is induced by $N_C$, so this is full faithfulness.
Every filtered colimit of representables is the nerve of the
corresponding colimit in $E$.  We obtain (2).  The preceding paragraphs
then prove (3), replacing $C$ by its essentially small retract closure.
Conversely, under (3), density identifies $X$ with its canonical colimit
of $E_\kappa/X$; this colimit exists by the assumed filteredness.
This is (1).

Finally let $E$ be cocomplete.  Colimits of fewer than $\kappa$
presentable objects are presentable: their Hom functors are limits of
fewer than $\kappa$ set-valued functors, and $\kappa$-filtered colimits
of sets commute with such limits.  The latter assertion follows by
choosing a common stage for fewer than $\kappa$ elements and a further
stage for fewer than $\kappa$ equalities.  Thus $E_\kappa/X$ is
$\kappa$-filtered whenever $E_\kappa$ is essentially small.
Starting only with a small dense $C$ of presentables, close it under
colimits of diagrams with fewer than $\kappa$ arrows.  Iterating for
$\kappa$ stages gives a small category $C'$ of presentables, closed
under these colimits.  It is still dense: a natural map on $C'$ is
uniquely determined on $C$, and the maps from $C$ jointly detect its
components.  The category $C'/X$ is filtered as above.  Applying the
already proved argument to $C'$ gives (1) and the required smallness
of all presentables.
\end{proof}

This is exactly the extra condition written on 161-2/060, conditions
(i)--(iii) of the converse.  The parenthetical uncertainty on scan 055
must not erase it.  The proof also spells out why the cocomplete case
on scans 061--062 is different.

\subsection{What the cardinal estimate actually says}
\begin{lemma}[Proved: small directed subdiagrams]\label{cat:I2}
Let $I$ be a $\kappa$-directed poset, let $\mu\geq\kappa$ be infinite,
and assume $\mu^{<\kappa}=\mu$.  Every subset of $I$ of size at most
$\mu$ is contained in a $\kappa$-directed subset of size at most $\mu$.
Let $J_\mu$ be the poset of these directed subsets, ordered by inclusion.
It is $\mu^+$-directed, its union is $I$, and
\[
 \operatorname{cf}(J_\mu)
 =\operatorname{cf}([I]^{\leq\mu},\subseteq).
\]
In particular, if $|I|\leq\nu$ and $\nu^\mu=\nu$, $J_\mu$ has at most
$\nu$ elements.  More generally a cofinal subposet with at most $\nu$
elements exists exactly when the displayed cofinality is at most $\nu$.
\end{lemma}
\begin{proof}
Starting with a subset $A_0$ of size at most $\mu$, at each successor
stage adjoin one chosen upper bound in $I$ for every subset of the
current set of size less than $\kappa$; take unions at limit stages.
There are at most $\mu^{<\kappa}=\mu$ choices at each stage.  After
$\kappa$ stages the resulting set still has size at most $\mu$ and is
$\kappa$-directed: every set of fewer than $\kappa$ of its elements
appears before one stage, by regularity.  The union of at most $\mu$
members of $J_\mu$ has size at most $\mu$, so the construction gives
an upper bound for it in $J_\mu$.  This proves directedness and also
that $J_\mu$ is cofinal in $[I]^{\leq\mu}$.  Cofinal subposets have
the same cofinality, which proves the formula.  The number of subsets
of size at most $\mu$ in a set of size at most $\nu$ is at most
$\nu^\mu$ (the finite exceptional cases are already excluded).
The last assertion is the definition of cofinality.
\end{proof}

\begin{proposition}[Proved: the comparison map in the cardinal filtration]\label{cat:I3}
Under the hypotheses of Lemma~\ref{cat:I2}, let $D:I\to E$ be a
diagram in a category having the relevant directed colimits.  Set
$X_J=\operatorname{colim}_{i\in J}D(i)$ for $J\in J_\mu$.
The canonical map
\[
 \operatorname{colim}_{J\in J_\mu}X_J
 \longrightarrow\operatorname{colim}_{i\in I}D(i)
\]
is an isomorphism.  Consequently the cardinal bound in the
manuscript's strengthened axiom c) is justified by the two closure
and indexing bounds of Lemma~\ref{cat:I2}.  An equality such as
$\nu^\mu=\nu$ is a sufficient bound, not a necessary condition for an
individual object to admit a small presentation.
\end{proposition}
\begin{proof}
A compatible collection of maps $X_J\to Y$ gives a map $D(i)\to Y$
by choosing any $J$ containing $i$.  Independence follows by choosing
an upper bound of two such $J$'s.  For $i\leq i'$, choose a $J$
containing both; the resulting maps form an $I$-cocone.  Conversely
an $I$-cocone restricts to every $J$ and hence gives the compatible
collection.  These constructions are mutually inverse and natural in
$Y$, proving the isomorphism by the universal property of the colimit.
The bounds are then precisely those in the preceding lemma.
\end{proof}

\begin{proposition}[Refuted: an accessible functor need not preserve all sufficiently large presentability ranks]\label{cat:I4}
The scan-055 assertion
$\exists\kappa_0\ \forall\kappa\geq\kappa_0:
F(E_\kappa)\subseteq F_\kappa$
is false even for an accessible endofunctor of $\mathbf{Set}$.
Take $F(X)=X^{\mathbf N}$.  This functor preserves
$\aleph_1$-filtered colimits, but fails to preserve
$\mu^+$-presentable objects for every infinite cardinal $\mu$ of
countable cofinality.  Such $\mu$ are arbitrarily large.
\end{proposition}
\begin{proof}
A function from $\mathbf N$ into a filtered colimit is represented at
one stage of an $\aleph_1$-filtered diagram, since its countably many
coordinates have a common stage.  Equality of two such functions is
witnessed at one further common stage.  This proves preservation.
In $\mathbf{Set}$ an object is $\lambda$-presentable, for regular
$\lambda$, exactly when its cardinality is less than $\lambda$.
One direction is the same element-and-equality argument; for the other,
write a set as the $\lambda$-directed union of its subsets of size
less than $\lambda$ and apply presentability to its identity.

Now $\mu^{\aleph_0}>\mu$.  For completeness, partition an alleged
list of at most $\mu$ functions $\mathbf N\to\mu$ into countably
many subsets $A_n$, each of cardinality less than $\mu$.  Choose the
$n$th coordinate of a new function outside the $n$th-coordinate
values occurring in $A_n$.  The new function is absent from the list.
Thus a set of size $\mu$ is $\mu^+$-presentable but its image under
$F$ is not.  Taking suprema of increasing countable sequences of
cardinals produces arbitrarily large such $\mu$.
\end{proof}

\begin{proposition}[Proved under stated hypotheses: reconstruction of functors]\label{cat:I5}
If $E$ and $F$ are $\kappa$-accessible and
$T:E\to F$ preserves $\kappa$-filtered colimits and sends
$\kappa$-presentables to $\kappa$-presentables, then the canonical
comparison
$\operatorname{Ind}_\kappa(T|_{E_\kappa})\simeq T$
is a natural isomorphism under Theorem~\ref{cat:I1}'s equivalences.
Conversely every functor obtained by this Ind construction has these
two properties.  The qualification on the rank is essential, as
Proposition~\ref{cat:I4} shows.
\end{proposition}
\begin{proof}
On a presentation $X=\operatorname{colim}_I x_i$ by presentables,
the comparison is the map
$\operatorname{colim}_I T(x_i)\to T(X)$ supplied by the cocone.
It is an isomorphism by preservation.  Using the canonical category
$E_\kappa/X$ makes these maps natural in $X$.  Conversely the Ind
extension is defined by these colimits and preserves them by the
same flattening of categories of elements as in
Theorem~\ref{cat:I1}.  A presentable object is a retract of a
representable, and its image is consequently a retract of a
presentable object of $F$, hence presentable.
\end{proof}

\subsection{The relative classifying topos}
All topoi in the following statement are Grothendieck topoi over the
fixed universe.  We use a consistent inverse-image convention for
2-morphisms.  Let $\mathbf B_I$ denote the classifier of a small
$I$-indexed family of objects.  For one sort it is
$[\mathbf{FinSet},\mathbf{Set}]$; for several sorts it is the
classifying topos of the empty theory with those sorts.  The notation
$\widehat{(\mathbf{FinSet})^{\mathrm{op}}}^{\,I}$ on scan 112 must
mean a product in the 2-category of topoi if it denotes this
classifier; the ordinary product of the underlying categories is
not the classifier of several independent objects.

\begin{theorem}[Proved, using the opened classifying-topos existence theorem]\label{cat:I6}
Let $T$ be a small geometric theory with classifying topos $B$, and
let $(b_i)_{i\in I}$ be specified sorts of its universal model.
For a Grothendieck topos $X$ and objects $E_i\in X$ there is a
Grothendieck topos $X[T;E]\to X$ such that, for every
$g:Y\to X$, geometric morphisms $Y\to X[T;E]$ over $X$ classify
$T$-models $M$ in $Y$ equipped with specified isomorphisms
$M_i\simeq g^*E_i$.
With $b:B\to\mathbf B_I$ and $e:X\to\mathbf B_I$ the morphisms
classifying the two families, we have a 2-pullback
\[
 X[T;E]\simeq B\times_{\mathbf B_I}X.
\]
This includes small finite-limit theories, as requested on
cote 161-2/112.
\end{theorem}
\begin{proof}
Choose a small generating full subcategory $C$ of $X$, enlarge it to
contain every $E_i$, and close it under finite limits.  This remains
small.  Equip $C$ with the induced canonical topology $J$; the
comparison functor gives $X\simeq\operatorname{Sh}(C,J)$.
The geometric theory $T_C$ has a sort for every object of $C$, a
function for every arrow, the equations for identities and composition,
the finite-limit axioms, and the following axiom for each covering
family $(c_a\to c)_a$:
\[
 x=x\ \vdash_{x:c}\ \bigvee_a\exists y:c_a\ (f_a(y)=x).
\]
Its models in $Y$ are exactly the left exact $J$-continuous functors
$C\to Y$, equivalently geometric morphisms $Y\to X$.
This identification, including morphisms, is the site-classification
result in Makkai--Reyes~\cite[Propositions 9.1.4--9.1.5]{catMakkaiReyes1977}.

Form the disjoint union of the signatures and axioms of $T_C$ and
$T$.  For every $i$, add two function symbols
$u_i:M_i\to F(E_i)$ and $v_i:F(E_i)\to M_i$, with equations
$v_i u_i=1$ and $u_i v_i=1$.  These are geometric axioms, so the
resulting small geometric theory $T'$ has a classifying topos by
\cite[Theorem 9.1.1]{catMakkaiReyes1977}.  Define
$X[T;E]=\operatorname{Set}[T']$.  Its universal $T_C$-model gives
its morphism to $X$.

By construction, a model of $T'$ in $Y$ is exactly a triple
$(g,M,(u_i)_i)$ as in the statement.  Morphisms of models are
exactly the compatible pairs of morphisms with the indicated
identifications, because the inverse equations impose no extra choices.
The equivalence is natural under inverse-image functors: those preserve
finite limits, arbitrary unions, and images, hence all the displayed
axioms.  The hom-category of triples is the pseudopullback of
$\operatorname{Hom}(Y,B)$ and $\operatorname{Hom}(Y,X)$ over
$\operatorname{Hom}(Y,\mathbf B_I)$.  This is precisely the
2-pullback universal property, and fixing $g$ gives the relative
universal property.  A finite-limit theory is geometric by expressing
its specified limits as existence-and-uniqueness axioms, so it is
included.
\end{proof}
